\section{Physics of the optimized solutions}
\label{sec:physics}

The potential lowers the volume-integrated exoticity by concentrating the
scalar variation into a narrower region near the throat. It does not weaken
the NEC violation everywhere. We examine this redistribution in three
selected $210\times42$ drainholes, denoted D0$^\star$, D21$^\star$ and
D30$^\star$. These are finer-grid selections, distinct from the search minima
in Tables~\ref{tab:models} and \ref{tab:X}; D21$^\star$ is the solution already
identified there. Each has its own massless reference at the same $R_e$, $J$
and asymmetry. In particular, both members of the stationary D21 comparison
have $J=\clmPhysicsDtwentyoneJ$.

\begin{figure*}[t]
\centering
\includegraphics[width=0.94\textwidth]{mechanism_d21.pdf}
\caption{Stationary profiles of D21$^\star$ and its massless reference at the same
$R_e$ and $J$. The signed equatorial proper distance $\ell$ is zero at the
minimum-area throat. Panels show the equatorial scalar gradient, the angularly
integrated contribution to $X$, the area radius of constant-$\eta$ surfaces,
and the energy of a zero-angular-momentum photon measured by a local observer
relative to its conserved energy at positive infinity. The potential
concentrates the scalar gradient and exoticity near the throat while
sharpening the interior redshift profile.}
\label{fig:mechanism}
\end{figure*}

\subsection{Matter distribution and throat geometry}

Figure~\ref{fig:mechanism} compares profiles against physical proper distance,
rather than the radial coordinates of different numerical solutions. The
local radial NEC deficit is the squared orthonormal scalar gradient divided
by $4\pi$. Here the radial null vector has unit energy in the local
zero-angular-momentum observer's frame. The potential cancels from its direct
null contraction, but changes the scalar and metric through the field equations.

In D21$^\star$, the scalar excursion between the two ends falls from
$\clmPhysicsExcursionClassical$ to $\clmPhysicsExcursionPotential$. Its remaining
variation occurs over a much shorter distance, so the peak radial NEC
violation is $\clmPhysicsPeakRatio$ times stronger. The proper width between
the 5\% and 95\% cumulative contributions to $X$ decreases from
$\clmPhysicsWidthClassical R_e$ to $\clmPhysicsWidthPotential R_e$. This measure
includes the long classical tails beyond the plotted interval; it describes
the matter distribution, not a change in throat radius. Angular scalar
gradients remain small, giving a narrow radial layer rather than an
equatorial ring.

The same pattern occurs in the other two solutions
(Table~\ref{tab:physics}). Their narrower distributions outweigh the higher
peaks, reducing $X$ by $\clmPhysicsReductionRange\%$ across the three models.
The equatorial throat radius is fixed, and the throat area changes by at most
$\clmPhysicsAreaChange\%$. The rotating throats become slightly less oblate.
Thus the reduction mainly reflects a changed scalar distribution, while the
overall throat dimensions remain similar.

\begin{table}[t]
\caption{Stationary quantities divided by their own massless-reference values
at fixed $R_e$, $J$ and asymmetry. Stars identify the selected finer-grid
solutions. The width includes the complete radial distribution; the null NEC
integral follows the complete equatorial ray with unit Killing energy at
positive infinity. Peak tides and light transit use equatorial proper
distances $-10R_e$ to $+10R_e$. A smaller volume integral accompanies stronger
local violation, stronger peak tides and longer transit.}
\label{tab:physics}
\small
\begin{ruledtabular}
\begin{tabular}{lrrr}
Quantity & D0$^\star$ & D21$^\star$ & D30$^\star$ \\
\hline
\clmPhysicsRows
\end{tabular}
\end{ruledtabular}
\end{table}

\subsection{Tides, redshift and light propagation}

A smaller volume integral does not imply a gentler passage through the
throat. We measure tides by the largest absolute eigenvalue of the electric
Riemann tensor in a radially moving observer's frame. At the same local
speed $v=0.5$, its peak over the matched interval increases in all three
solutions. This compares local forces along a fixed-speed scan; a freely
falling traveller's accumulated tidal exposure also depends on its trajectory.
The stronger peak is consistent with the sharper scalar and metric variation
visible in Fig.~\ref{fig:mechanism}.

The asymptotic redshift remains fixed by the common asymmetry $\gamma$,
while the interior profile changes. For a zero-angular-momentum photon,
the local energy is $E_{\mathrm{local}}/E_{+\infty}=1/\alpha$, where
$\alpha=e^{f/2}$ is the lapse and $E_{+\infty}$ is the conserved energy
normalized at positive infinity. The optimized profile makes the transition
between the two ends more abrupt.

We compute a finite light-crossing time along the equatorial null geodesic
by integrating $dt/d\ell=1/\alpha$ between endpoints at equal proper
distances from each throat. The same positive-infinity clock normalizes
$t$ in both solutions. Transit is longer for the optimized solutions
(Table~\ref{tab:physics}); the increase also occurs when the endpoint
distances are halved or doubled, although the ratio changes. These are
wormhole-to-wormhole comparisons at specified endpoints.

The magnitude of the NEC integral along the complete equatorial null ray
also increases. Its affine parameter is normalized by unit Killing energy
at positive infinity, unlike the local normalization used for the peak NEC
deficit. The ray and volume integrals therefore measure different aspects
of the same distribution. Minimizing $X$ reduces the volume-weighted matter
requirement while concentrating the violation encountered near the throat.
